\section{Regarding Bayes Factors}\label{app:bayes_factors}

This appendix concerns the use of the Bayes Factor (BF) as a \textit{sequential
stopping rule}; this role is distinct from its use in Bayes Factor Design Analysis
(BFDA), which is a prospective planning tool discussed in Section~\ref{sec:discussion}.

The BF is a widely used Bayesian measure for comparing two hypotheses.
For a null hypothesis $H_0$ and an alternative $H_1$, the Bayes Factor in favour of
the null is defined as the ratio of marginal likelihoods:
\begin{equation}\label{eq:bf_definition}
  \mathrm{BF}_{01} = \frac{p(\mathrm{data} \mid H_0)}{p(\mathrm{data} \mid H_1)}.
\end{equation}
A value $\mathrm{BF}_{01} > 1$ constitutes evidence in favour of $H_0$; a value
$\mathrm{BF}_{01} < 1$ constitutes evidence against it.
The scale proposed by \citet{jeffreys1961} and refined by \citet{kassraftery1995}
provides conventional thresholds: $\mathrm{BF}_{01} > 3$ is typically taken as
``substantial'' evidence for the null, and $\mathrm{BF}_{01} < 1/3$ as substantial
evidence against it.

For the Bernoulli setting with $z$ successes in $N$ trials, \citet{kruschke2015doing}
(Chapter~13.3.2) uses a point null $H_0\colon \theta = \theta_{\rm null}$ and a uniform
alternative $H_1\colon \theta \sim \mathrm{Beta}(1, 1)$.  The marginal likelihood under
$H_0$ is simply $\theta_{\rm null}^{z}(1-\theta_{\rm null})^{N-z}$, and under $H_1$ it
is the Beta--Binomial integral
\begin{equation}\label{eq:bf_h1_marginal}
  p(\mathrm{data} \mid H_1)
  = \int_0^1 \theta^{z}(1-\theta)^{N-z}\,d\theta
  = B(z+1,\, N-z+1),
\end{equation}
where $B(\cdot,\cdot)$ denotes the Beta function, giving the closed-form expression
\begin{equation}\label{eq:bf_bernoulli}
  \mathrm{BF}_{01}
  = \frac{\theta_{\rm null}^{z}(1-\theta_{\rm null})^{N-z}}{B(z+1,\, N-z+1)}.
\end{equation}
The sequential stopping rule is: continue sampling until either
$\mathrm{BF}_{01} \geq 3$ (accept $H_0$) or $\mathrm{BF}_{01} \leq 1/3$ (reject $H_0$),
or the budget $N_{\rm max}$ is exhausted.

\subsection*{Sequential Behaviour}

\citet{kruschke2015doing} evaluates the operating characteristics of BF stopping
sequentially across 1{,}000 simulated sequences of up to $N_{\rm max}=1{,}500$ flips,
comparing it against NHST, HDI+ROPE, and PitG (his Figures~13.6 and~13.7).
The main findings relevant to this paper are as follows.

\subsubsection*{When the null is true ($\theta_{\rm true}=\theta_{\rm null}=0.5$)}

This is the setting shown in Kruschke's Figure~13.6 (for which
Figure~\ref{fig:fair_decisions} of this paper is analogous).
The BF stopping rule converges to an asymptotic false rejection rate
(referred to by \citeauthor{kruschke2015doing} as the \textit{false alarm rate})
of just over 20\%---lower than NHST, which diverges to 100\%, but substantially higher
than PitG, which maintains a 0\% false rejection rate \citep{kruschke2015doing}.
As demonstrated in this paper, DPitG also achieves a 0\% false rejection rate.
The BF correctly accepts the null for the remaining sequences, although it can do
so on the basis of precise or imprecise data indiscriminately.

\subsubsection*{When the null is false ($\theta_{\rm true}=0.65$)}

This setting is displayed in Kruschke's Figure~13.7 (which is captured in the middle
panel of Figure~\ref{fig:conclusiveness_rates}).
The BF falsely accepts the null in approximately 40\% of sequences despite the true
parameter lying outside the ROPE, a substantially higher false acceptance rate than
either precision-based method.

\subsubsection*{Parameter estimation at the stopping point}

Because the BF stopping rule fires on the \emph{extremeness} of the likelihood ratio
rather than the \emph{width} of the posterior, it can accept or reject with very little
data.  As Kruschke's lower rows illustrate, $\hat{\theta}$ at stopping is noticeably
biased away from $\theta_{\rm true}$ in both the accept and the reject case---more so
than HDI+ROPE, and substantially more than PitG.  This is the fundamental problem
identified in \citet{kruschke2015doing}: a stopping rule based on extremeness in the data
automatically biases the sample toward extreme estimates; only a precision-based stopping
rule guarantees near-unbiased parameter estimation.

\subsection*{Key Limitations in Context}

Three structural limitations follow from the BF's definition.

\subsubsection*{No effect-size criterion}

$\mathrm{BF}_{01}$ compares the point null $\theta = \theta_{\rm null}$ against the
prior-averaged alternative.  It carries no information about whether the true parameter
is \emph{practically} different from the null---a parameter value just outside the ROPE
is treated identically to one that is far from it.  The HDI+ROPE framework addresses
this explicitly by requiring the HDI to fall cleanly inside or outside the ROPE before
a verdict is issued.

\subsubsection*{Sensitivity to the alternative prior}

$\mathrm{BF}_{01}$ depends on the choice of prior $p(\theta \mid H_1)$
(Equation~\ref{eq:bf_h1_marginal}).  Kruschke's choice of a uniform
$\mathrm{Beta}(1,1)$ prior is one convention; different priors yield different
numerical BF values and therefore different stopping points, even for identical data.
The HDI+ROPE stopping rule has no such ambiguity: its behaviour depends only on the
posterior and the pre-specified ROPE and $\omega_{\rm goal}$.

\subsubsection*{No precision guarantee}

The BF can stop early on imprecise posteriors.  In contrast, PitG and DPitG require
the HDI width to fall below $\omega_{\rm goal}$ before any verdict is issued, providing
a structural guarantee that accepted or rejected conclusions rest on adequately precise
estimates.

These limitations notwithstanding, BF stopping and HDI+ROPE are not mutually exclusive.
As noted in Section~\ref{sec:discussion}, a design requiring both a conclusive
HDI+ROPE verdict \textit{and} a minimum Bayes Factor is a natural extension of DPitG,
combining the precision guarantee of the latter with the evidential strength measure
of the former.
